The multi-agent pipeline developed in this work emulates and extends beyond human processes of hypothesis generation by synthesizing information around a design problem and extrapolating to adjacent domains of knowledge. We make use of targeted injections within in-context learning of LLMs which can critically shape model behavior and guide generated output. In this method, an injection refers to the deliberate inclusion of structured knowledge, constraints, or illustrative examples into the model’s immediate context, thereby ensuring that its reasoning process remains aligned with user objectives. By curating the information presented during inference, researchers can direct the model toward recognizing patterns or relationships outside the LLMs training data that might otherwise remain hidden. This approach is particularly valuable in complex domains such as materials discovery, where the design space is vast and unconstrained. Well-designed injections allow the model to focus on salient properties, standards, and experimental protocols, thereby improving both the reliability and interpretability of its outputs. Through this mechanism, in-context injections transform the model from a generic language processor into a specialized assistant capable of supporting scientific problem solving.

We demonstrate this approach in the search for sustainable alternatives to PFAS materials, where in-context content is dynamically tuned using external knowledge bases and tools accessible to agents on demand. In particular, we integrate knowledge graphs of PFAS and their intrinsic material properties, enabling agents to navigate a broader knowledge space and apply graph traversal algorithms to uncover latent connections embedded in the graph but not yet explicitly recognized in the literature. We stipulate that the materials and design hypotheses discussed in the PFAS use case are presented to illustrate the behavior and capabilities of the multi-agent system and to highlight early-stage development potential, rather than to prescribe implementation-ready material substitutions. It is important to note that the study and examples are used to illustrate the models rather than being final solutions, where further work is necessary, including experimental validation.

By design, our knowledge graphs currently generate hypotheses rooted in an academic corpus, leading to proposed composites that emphasize scientifically novel or exotic materials (e.g., graphene oxide, cellulose nanocrystals, lignin blends) but are not always industrially viable. To overcome this limitation, future work will expand the corpus to include diverse sources such as patents and technical datasheets, which capture industry-relevant testing protocols, performance standards, and material metrics. Incorporating these sources will enhance design-problem formulation and facilitate the identification of key design parameters, enabling tighter alignment with corresponding nodes in the global knowledge graph. This expansion will, in turn, improve the accuracy and realizability of the predicted property profiles, ensuring closer conformity with the design objectives and performance characteristics of PFAS. Further, wider integration of external datasets will broaden the range of traversable paths within the graph, enriching the substrate from which the Engineer agent can derive candidate designs. Providing the agentic system with greater access to validation tools would establish stronger feedback loops, improving the generated hypotheses and guiding iterative refinement of material designs. 

Currently, GraphAgents relies on pre-computed, static knowledge graphs. A logical next step is to enable agents to construct their own reasoning substrates dynamically during inference. Future iterations could integrate \textit{in-situ} knowledge expansion mechanisms, such as the Graph-PReFLexOR framework~\cite{Buehler2025GraphPRefLexOR,Buehler2025Chaos}, allowing the system to build temporary `scaffolding' graphs to bridge distant domains only when necessary. Furthermore, coupling this creative engine with autonomous synthesis labs would close the loop, transforming these digital hypotheses into physical reality without human intervention.